\vspace{-2mm}
\section{Conclusion and Future Work}
\label{conclusion}
Eight-bit weight--activation \gls{QAT} preserves link performance and
reduces modeled feature-map traffic in a standard-compliant 5G~NR
MU-\gls{MIMO} neural receiver. FP4 outperforms \gls{LS}--\gls{LMMSE}, whereas INT4
performs poorly, supporting microscaling formats for neural receivers \cite{rouhani2023microscaling,ocp2023mx}. Analytically, the weight-activation quantization followed by pruning achieves $66\times$ and
$8.8\times$ reductions in bit-operations and weight storage,
respectively. This robustness aligns with FP4's denser near-zero grid
and lower residual-path sparsity.

Future work will pursue hardware-efficient 2:4 structured sparsity and native
MXFP4 block scaling. Distillation-augmented \gls{QAT} offers a
complementary route to further improve ultra-low-bit neural receiver link performance
\cite{yawei_qat_distillation}.